\section{Online Evaluation}
\label{sec:eval}

We evaluate the savings available under PACE's cumulative budget, its comparison with published decision rules, and the contribution of each component.
The study uses measured cost profiles on synthetic streams and recorded arrival orders.

\paragraph{Setup.}
We run two grids of $300$ conditions, each crossing three models, four verification settings, five arrival patterns, and five cost or environment settings, with eight paired repetitions per condition.
The second grid changes seeds and parameters after the first study.
We also use Sepsis and BPI 2019 process logs \citep{mannhardt2016sepsis,vandongen2019bpi} and two Wikipedia edit streams, giving twelve combinations of models and logs with ten paired repetitions.
The logs supply task repetitions and their order, not live GUI tasks.
We shuffle measured cost profiles onto families and assign verification probabilities independently of initial compilation outcomes.
Unobserved successful or failed costs require substitutions, whose effect we test below.
The policies receive family labels and failure parameters, including detected-break status.
Appendix~\ref{app:simulation-details} specifies the grids, mappings, substitutions, and shared outcome sequences.

We compare PACE with ReAct \citep{yao2023react} and decision-rule adaptations of AutoRPA \citep{chen2026autorpa} and ToolPro \citep{liu2026toolpro} under the same compiler and cost model.
ReAct serves every arrival with the agent.
AutoRPA makes one compilation attempt when enough source traces are available and does not retry a rejected attempt.
ToolPro compares the compilation cost with the saving available from the current task, without projecting later uses.
These adaptations isolate when to invoke the shared compiler.
All policies pay the same routing, extraction, fallback, and realized compilation charges, so the comparison does not attribute differences in compiler quality to the decision rule.
Each table ratio divides a policy's mean cost by PACE's mean cost in the same condition.
Aggregate results are the equally weighted means of these ratios.
Reported percentage savings are the equally weighted means of $100(1-K_{\mathrm{PACE}}/K_{\mathrm{baseline}})$ across conditions, using each condition's mean costs.
We use $2000$ paired repetition-bootstrap draws, conditional on the tested conditions and supplied profiles.
These are development comparisons rather than independent confirmation after protocol selection.
The main tables give point estimates, with complete conditional intervals in Appendix~\ref{app:full-comparisons}.

\subsection{Savings}
\label{sec:regimes}
\label{sec:realstreams}

Across the grids, ReAct costs $1.251$ and $1.289$ times PACE on average (Table~\ref{tab:policysim}).
No saved prefix exceeds the implemented $1.25$ bound over $4800$ grid runs and $120$ log runs.
This checks the implementation on those runs.
Theorem~\ref{thm:cost-protection} covers all histories allowed by the model.
AutoRPA's largest grid ratios to PACE are $14.085$ and $19.269$, while ToolPro stays near ReAct cost by allowing no reuse beyond one task in its compilation criterion.

% Generated by rewrite-review/render_pace_tables.py from frozen evidence.
\begin{table}[!htbp]
\centering
\caption{Cost ratios relative to PACE over two grids of 300 conditions each.
Mean is the mean of the ratios for individual conditions, and max is the largest tested condition mean.
Blue, amber, and green identify GLM, DeepSeek, and Qwen, respectively.
Bold marks the best displayed values.
Conditional 95\% intervals are in Appendix~\ref{app:full-comparisons}.}
\label{tab:policysim}
\tablefont
\setlength{\tabcolsep}{0.6pt}
\renewcommand{\arraystretch}{0.95}
\begin{tabular}{@{}lrrrr@{}}
\toprule
\multicolumn{1}{c}{Policy} & \multicolumn{1}{c}{\shortstack{Grid 1 mean\\($\times$ PACE) $\downarrow$}} & \multicolumn{1}{c}{\shortstack{Grid 1 max\\($\times$ PACE) $\downarrow$}} & \multicolumn{1}{c}{\shortstack{Grid 2 mean\\($\times$ PACE) $\downarrow$}} & \multicolumn{1}{c}{\shortstack{Grid 2 max\\($\times$ PACE) $\downarrow$}} \\
\midrule
ReAct & \modelcell{3.4em}{1.528}{1.155}{1.069} & \modelcell{3.4em}{4.543}{4.070}{2.232} & \modelcell{3.4em}{1.585}{1.190}{1.090} & \modelcell{3.4em}{5.487}{3.247}{2.979} \\
AutoRPA & \modelcell{3.4em}{1.542}{2.226}{2.375} & \modelcell{3.4em}{4.234}{13.875}{14.085} & \modelcell{3.4em}{1.601}{2.296}{2.471} & \modelcell{3.4em}{5.178}{18.652}{19.269} \\
ToolPro & \modelcell{3.4em}{1.506}{1.155}{1.069} & \modelcell{3.4em}{4.431}{4.070}{2.232} & \modelcell{3.4em}{1.554}{1.190}{1.090} & \modelcell{3.4em}{5.198}{3.247}{2.979} \\
\midrule
\textbf{PACE} & \modelcell{3.4em}{\textbf{1.000}}{\textbf{1.000}}{\textbf{1.000}} & \modelcell{3.4em}{\textbf{1.000}}{\textbf{1.000}}{\textbf{1.000}} & \modelcell{3.4em}{\textbf{1.000}}{\textbf{1.000}}{\textbf{1.000}} & \modelcell{3.4em}{\textbf{1.000}}{\textbf{1.000}}{\textbf{1.000}} \\
\bottomrule
\end{tabular}
\end{table}

\begin{table}[!htbp]
\centering
\caption{Cost ratios relative to PACE on four recorded arrival sequences, using ten paired repetitions.
Blue, amber, and green identify GLM, DeepSeek, and Qwen, respectively.
Bold marks the best displayed values.
Conditional 95\% intervals are in Appendix~\ref{app:full-comparisons}.}
\label{tab:realstreams}
\tablefont
\setlength{\tabcolsep}{0.6pt}
\renewcommand{\arraystretch}{0.95}
\begin{tabular}{@{}lrrrr@{}}
\toprule
\multicolumn{1}{c}{Policy} & \multicolumn{1}{c}{\shortstack{Sepsis\\($\times$ PACE) $\downarrow$}} & \multicolumn{1}{c}{\shortstack{BPI 2019\\($\times$ PACE) $\downarrow$}} & \multicolumn{1}{c}{\shortstack{Wiki tools\\($\times$ PACE) $\downarrow$}} & \multicolumn{1}{c}{\shortstack{Wiki humans\\($\times$ PACE) $\downarrow$}} \\
\midrule
ReAct & \modelcell{3.4em}{0.972}{1.011}{\textbf{0.993}} & \modelcell{3.4em}{2.236}{1.571}{1.480} & \modelcell{3.4em}{2.511}{1.510}{1.381} & \modelcell{3.4em}{1.112}{1.056}{1.039} \\
AutoRPA & \modelcell{3.4em}{0.984}{1.170}{1.219} & \modelcell{3.4em}{2.198}{1.628}{1.559} & \modelcell{3.4em}{2.462}{1.542}{1.416} & \modelcell{3.4em}{1.095}{1.168}{1.197} \\
ToolPro & \modelcell{3.4em}{\textbf{0.968}}{1.011}{\textbf{0.993}} & \modelcell{3.4em}{2.154}{1.571}{1.480} & \modelcell{3.4em}{2.482}{1.510}{1.381} & \modelcell{3.4em}{1.105}{1.056}{1.039} \\
\midrule
\textbf{PACE} & \modelcell{3.4em}{1.000}{\textbf{1.000}}{1.000} & \modelcell{3.4em}{\textbf{1.000}}{\textbf{1.000}}{\textbf{1.000}} & \modelcell{3.4em}{\textbf{1.000}}{\textbf{1.000}}{\textbf{1.000}} & \modelcell{3.4em}{\textbf{1.000}}{\textbf{1.000}}{\textbf{1.000}} \\
\bottomrule
\end{tabular}
\end{table}


On the recorded logs, PACE reduces token costs by $21.85\%$ relative to ReAct, with a conditional $95\%$ interval of $[19.66,24.15]\%$ (Table~\ref{tab:realstreams}).
It saves cost in ten of twelve conditions.
Its largest cost ratio to ReAct is $1.029$, and its lowest is $0.398$.
Relative to PACE, AutoRPA and ToolPro have mean cost ratios of $1.470$ and $1.396$, respectively, although they cost less than PACE in one and two conditions.
The comparison concerns the stated decision-rule adaptations, rather than the original systems end to end.

\subsection{Component effects}
\label{sec:ablation}

Table~\ref{tab:ablation} reports cost ratios relative to full PACE.
\emph{No savings test} compiles whenever eligible while retaining the budget.
\emph{No cost budget} removes reservations and budget-triggered manifest clearing.
\emph{Eager retry} removes both, and \emph{Equal attempt costs} substitutes $C/\hat p$ for $\hat B$ while retaining actual charges and reservations.

% Generated by rewrite-review/render_pace_tables.py from frozen evidence.
\begin{table}[htbp]
\centering
\caption{Cost ratios after removing or simplifying PACE components.
Values are mean cost ratios relative to PACE, which is exactly one.
Blue, amber, and green identify GLM, DeepSeek, and Qwen, respectively.
Paired intervals are in Appendix~\ref{app:ablation-absolute}.
Grid max is the largest tested condition-mean ratio, without a sampling interval.}
\label{tab:ablation}
\tablefont
\setlength{\tabcolsep}{0.9pt}
\renewcommand{\arraystretch}{0.95}
\begin{tabular}{@{}lrrrr@{}}
\toprule
\multicolumn{1}{c}{Variant} & \multicolumn{1}{c}{\shortstack{Grid 1 mean\\($\times$ PACE)}} & \multicolumn{1}{c}{\shortstack{Grid 2 mean\\($\times$ PACE)}} & \multicolumn{1}{c}{\shortstack{Logs mean\\($\times$ PACE)}} & \multicolumn{1}{c}{\shortstack{Grid max\\($\times$ PACE)}} \\
\midrule
\textbf{PACE} & \modelcell{3.4em}{1.000}{1.000}{1.000} & \modelcell{3.4em}{1.000}{1.000}{1.000} & \modelcell{3.4em}{1.000}{1.000}{1.000} & \modelcell{3.4em}{1.000}{1.000}{1.000} \\
No savings test & \modelcell{3.4em}{1.104}{1.185}{1.204} & \modelcell{3.4em}{1.106}{1.144}{1.166} & \modelcell{3.4em}{1.327}{1.568}{1.499} & \modelcell{3.4em}{1.881}{1.859}{1.788} \\
No cost budget & \modelcell{3.4em}{1.028}{1.015}{1.007} & \modelcell{3.4em}{1.028}{1.016}{1.012} & \modelcell{3.4em}{1.000}{1.000}{1.000} & \modelcell{3.4em}{1.472}{1.397}{1.302} \\
Equal attempt costs & \modelcell{3.4em}{1.008}{1.115}{1.083} & \modelcell{3.4em}{1.007}{1.088}{1.071} & \modelcell{3.4em}{1.000}{1.312}{1.088} & \modelcell{3.4em}{1.402}{1.585}{1.319} \\
Eager retry & \modelcell{3.4em}{2.321}{11.756}{12.650} & \modelcell{3.4em}{2.385}{13.322}{14.353} & \modelcell{3.4em}{1.395}{4.487}{3.930} & \modelcell{3.4em}{14.813}{111.240}{151.251} \\
\bottomrule
\end{tabular}
\end{table}


Removing the savings test raises paired mean log cost by $46.46\%$ (conditional $95\%$ interval $[41.60,51.74]\%$).
Ignoring failed compilation costs raises it by $13.36\%$ ($[12.05,15.04]\%$).
The GLM log effect of equal costs is exactly zero because those profiles already set $C^{\mathrm{fail}}=C$.
Both variants keep the budget, so their largest grid ratios to ReAct remain near the same $1.25$ limit despite higher mean costs.
Removing the budget changes mean log cost by only $-0.014\%$, yet raises the largest tested grid ratio to ReAct from $1.2494$ to $1.8387$, or $47.2\%$.
Thus the near-zero log effect concerns average cost, while the budget's measured contribution is protection against larger losses.
Appendix~\ref{app:ablation-absolute} gives the cost ratios and paired intervals, including conditions where an ablation is cheaper.

\paragraph{Cost assumptions.}
Varying the missing failed/successful compilation cost ratio over $0.5,1,2,5$ gives PACE mean log ratios from $0.788$ to $0.830$.
Assigning high compilation costs to frequent families with low verification probability raises the mean to $0.997$, nearly removing the saving while preserving the cost constraint.
This assignment makes the families with the most reuse opportunities costly to acquire.
The resulting loss of savings shows why task frequency alone is insufficient to select compilations.
The budget still limits extra spending, but it cannot create savings when the available compilation costs and useful lifetimes make reuse unprofitable.
Appendix~\ref{app:cost-sensitivity} retains all seventy-two sensitivity conditions.
Appendix~\ref{app:live} reports a separate live GUI component study.
